% =====================================================================
\chapter{Perspectivas futuras}\label{cap:perspectivas}
% Plantilla ITBA: opcional, máximo 1 carilla. Puede incluir objetivos de
% máxima no alcanzados. Borrador del 21 oct. 2026: incluye, como trabajo a
% futuro, las limitaciones del diseño que la Dra. Kruse sacó de la Discusión
% (a confirmar con ella).
% =====================================================================

Los resultados de este trabajo abren distintas líneas para continuarlo:

\begin{itemize}
  \item \textbf{Neurogénesis en los grupos expuestos a stevia.} Cuantificar PCNA y DCX en los animales expuestos a stevia permitirá determinar si la preservación de la separación de patrones espaciales se acompaña de una neurogénesis conservada. Conviene, además, registrar el animal de origen de cada corte, de modo que el análisis pueda realizarse por animal y no por corte.
  \item \textbf{Mecanismos.} Evaluar si las modificaciones en la expresión de calretinina y parvalbúmina descriptas tras el consumo agudo de stevia persisten a largo plazo y se relacionan con las alteraciones de la memoria de reconocimiento, y estudiar el papel de la abstinencia de la bebida endulzada en los cambios de la conducta tipo ansiosa.
  \item \textbf{Hembras.} Extender el estudio a ratas hembra, en las que el mismo extracto redujo la conducta tipo ansiosa y alteró la fertilidad, para establecer si esos efectos persisten a largo plazo y si dependen de la edad de exposición.
  \item \textbf{Dosis.} Dado que la ingesta superó los valores diarios de referencia, evaluar concentraciones menores del extracto, cercanas a la ingesta diaria admisible de glucósidos de esteviol.
  \item \textbf{Diseño de las cohortes.} Incluir un grupo control en cada cohorte experimental y ampliar el número de animales expuestos a stevia en el \textit{open field}.
  \item \textbf{Análisis multivariado.} Registrar las métricas de \anymaze{} identificadas por animal también en el \textit{open field} y en el NOR, para extender el PCA, el clustering y el Random Forest a toda la batería conductual e integrar en un mismo análisis las distintas pruebas, junto con el peso corporal.
\end{itemize}
